\begin{acknowledgements}
\begin{spacing}{1.32}

First and foremost, I am deeply grateful to my supervisors, Prof.~Ping Luo and Prof.~Hongyang Li. Prof.~Luo gave me the freedom to pursue the questions I cared about and the perspective to see how they fit together. Prof.~Li introduced me to autonomous driving research at OpenDriveLab and taught me, by example, that a good benchmark can move an entire community. This thesis would not exist without their guidance, patience and trust.

I thank Dr.~Kashyap Chitta and Prof.~Andreas Geiger, with whom I worked closely on DriveLM and Centaur; their rigour in evaluation and clarity in writing have shaped how I do research. I thank Dr.~Jos\'e M. \'Alvarez and Dr.~Zhiding Yu for mentoring me during my internship at NVIDIA Autonomous Vehicle Applied Research, where the ideas behind Centaur took shape. I am also grateful to Prof.~Yexiang Xue, who advised me during my years at Purdue University, and to Prof.~Kun He, who first guided me into research at Huazhong University of Science and Technology.

In the interest of research integrity, I wish to acknowledge my co-authors on the works forming this dissertation. Much of this research was conducted together with colleagues at OpenDriveLab, Shanghai AI Laboratory, and MMLab@HKU. Dr.~Li Chen has been an indispensable collaborator across all four core chapters, and Wenwen Tong, Tai Wang, Katrin Renz, Dr.~Kashyap Chitta, Checheng Yu, Modi Shi, Lirui Zhao, and Qingwen Bu made foundational contributions to the respective chapters of this thesis. I also thank Silei Wu, Hanming Deng, Yi Gu, Lewei Lu, Prof.~Dahua Lin, Hanxue Zhang, Chengen Xie, Jens Bei{\ss}wenger, Shiyi Lan, Yibo Yuan, Tianyu Li, Jin Chen, Kaiyang Wu, Shijia Peng, Huijie Wang, Haoguang Mai, Gangcheng Jiang, Hai Zhang, Jiazhi Yang, Yihan Hu, and Zhiqi Li for their collaboration, as well as the organisers, participants, and sponsors of the community challenges we hosted.

Parts of the research in this thesis were supported by the General Research Fund of Hong Kong (Nos.~17200622 and 17209324) and the National Key R\&D Program of China (No.~2022ZD0161000), as acknowledged in the corresponding publications. During my doctoral study, I was supported by a Senior Research Assistant position.

Finally, I thank my family for their unconditional love and support throughout this long journey.

\end{spacing}
\end{acknowledgements}
